% !PRINTED 248
\weberpage{248}
If, then, any element $\alpha$ belongs to the exponent $e$, the number of such elements equals the number of positive integers less than $e$ and relatively prime to $e$. This number, which we have already denoted by $\varphi(e)$, satisfies, as $e$ runs through all divisors of $p^n - 1$, the relation

\weq{6}{\textstyle\sum \varphi(e) = p^n - 1}

(Bd.\,I, \S.\,132, 133).

If $\psi(e)$ is the number of elements $\alpha$ that belong to the exponent $e$, then $\psi(e) = \varphi(e)$ or $= 0$, and since the total number of elements $\alpha$ is $p^n - 1$ and every element $\alpha$ belongs to one and only one exponent, we have also

\weq{7}{\textstyle\sum \psi(e) = p^n - 1,}

and from (6) and (7) it follows that $\psi(e)$ is \gesperrt{always equal to} $\varphi(e)$.

There are therefore $\varphi(p^n - 1)$ numbers $\gamma$ in $\mathfrak{C}$ that belong to the exponent $p^n - 1$, and which consequently have the property that the powers

\weq{8}{1, \gamma, \gamma^2, \ldots, \gamma^{p^n-2}}

are all different from one another. By this sequence, therefore, the totality of the numbers in $\mathfrak{C}$ different from zero is exhausted.

Such numbers $\gamma$ are called \gesperrt{primitive roots of $\mathfrak{C}$}.

Among the numbers in $\mathfrak{C}$ there are some that can be represented as the square of another number in $\mathfrak{C}$, which we call \gesperrt{squares}, and others for which this is not the case, which we call \gesperrt{non-squares}.

When $p = 2$, every number in $\mathfrak{C}$ is a square, as formula (4) shows.

But when $p$ is odd, one half of the numbers in $\mathfrak{C}$ different from zero consists of squares, the other of non-squares.

For two opposite numbers, such as $+\alpha$ and $-\alpha$, are different from one another when $p$ is odd, and yet give the same square. If, then, one raises all the numbers of $\mathfrak{C}$ to the square, one obtains at most $\frac{1}{2}(p^n - 1)$ different squares. Now, on the other hand, $\beta^2$ can be equal to $\alpha^2$ only when $\beta = \pm\alpha$; for from $\beta^2 - \alpha^2 = (\beta - \alpha)(\beta + \alpha) = 0$ it follows that $\beta - \alpha$ or $\beta + \alpha$ must vanish. There are therefore really $\frac{1}{2}(p^n - 1)$ squares and just as many non-squares.
